\section{Chapter 9}

\begin{verse}
To preserve when one is filled up to the brim\\
Is not possible\\
To keep a blade very sharp\\
Is not possible.\\
One cannot, at the same time,\\
Possess and conserve.\\
Goods and power united to pride\\
Prepare your own ruin\\
To act and pull back = to pass into the shadow \\
Is the Way of Heaven
\end{verse}

References, through images, to the law of converting everything that is joined at the extreme into the opposite. Decadence begins at the peak. For essential mastery and for control of the transformations, it is necessary to know how to abandon at the right point what is brought to completion. To be tied to, and especially to be associated with, the pride of the I, means to plummet. The inner meaning of these maxims must be placed before the social and political meaning, where it usually ended in banal terms (for example: not to accumulate riches and goods because they cannot be made secure --- or rather: to retreat to an obscure life after having realized or gathered honors).

\paragraph{Chinese text and literal translation}


\textbf{Chapter 9 (\begin{CJK*}{UTF8}{bsmi}第九章\end{CJK*})}
\columnratio{.4}
\begin{paracol}{2}
\begin{verse}
\begin{CJK*}{UTF8}{bsmi}
持而盈之，\\
不如其已。 \\
揣而銳之，不可常保。\\
金玉滿堂，莫之能守。\\
富貴而驕，自遺其咎。\\
功遂身退，天之道。
\end{CJK*}
\end{verse}
\switchcolumn
\begin{verse}
Rather than pour into a cup overflowed,\\
It is better to stop it.\\
To hammer and to sharpen, it will not be long to be broken.\\
One cannot keep in eternal hiding a hall brimming with riches.\\
The man arrogant from enormous wealth buries disaster upon his road.\\
Removing after succeeding fits in with the Tao/principle of the universe.
\end{verse}
\end{paracol}


\flrightit{Posted on 2020-09-27 by Lao Tzu}

